\begin{abstract}

Zero-shot time-series foundation models are usually reached by pretraining on corpora of tens or
hundreds of billions of observations. We ask how far a competitive zero-shot forecaster can be pushed
when that budget is small, by starting from a backbone already pretrained on temporal data of another
kind: video. The transfer does not work out of the box --- continuing a video model's own
masked-pixel objective on rendered series improves that objective while making forecasts steadily
worse --- so we build the adaptation stage around the forecasting task instead. A series is rendered
one period per frame, so that motion across frames is evolution from one period to the next; a
trailing block of frames is predicted from the visible past alone; and a differentiable read-out
decodes the predicted frames back into values, which is where the loss is taken. Training samples
the same frame scales the inference rule uses. After 60k optimizer updates at an effective batch of
32 on a corpus of 654M observations, roughly $0.3\%$ the size of the archive the strongest visual
baseline trains on, one checkpoint reaches a mean MSE of $0.285$ across six standard long-sequence
datasets and four horizons --- the lowest among the models we compare --- though its mean MAE,
$0.328$, is behind the best baseline's $0.326$. A controlled study that varies only the
initialisation shows the video prior is worth $3$--$5\%$ over an ImageNet prior at a 5k-update
budget and that the gap closes by 20k, so what video pretraining buys here is adaptation efficiency
rather than a higher ceiling.

\end{abstract}
